% Options for packages loaded elsewhere
\PassOptionsToPackage{unicode}{hyperref}
\PassOptionsToPackage{hyphens}{url}
\documentclass[
]{article}
\usepackage{xcolor}
\usepackage[margin=1in]{geometry}
\usepackage{amsmath,amssymb}
\setcounter{secnumdepth}{5}
\usepackage{iftex}
\ifPDFTeX
  \usepackage[T1]{fontenc}
  \usepackage[utf8]{inputenc}
  \usepackage{textcomp} % provide euro and other symbols
\else % if luatex or xetex
  \usepackage{unicode-math} % this also loads fontspec
  \defaultfontfeatures{Scale=MatchLowercase}
  \defaultfontfeatures[\rmfamily]{Ligatures=TeX,Scale=1}
\fi
\usepackage{lmodern}
\ifPDFTeX\else
  % xetex/luatex font selection
\fi
% Use upquote if available, for straight quotes in verbatim environments
\IfFileExists{upquote.sty}{\usepackage{upquote}}{}
\IfFileExists{microtype.sty}{% use microtype if available
  \usepackage[]{microtype}
  \UseMicrotypeSet[protrusion]{basicmath} % disable protrusion for tt fonts
}{}
\makeatletter
\@ifundefined{KOMAClassName}{% if non-KOMA class
  \IfFileExists{parskip.sty}{%
    \usepackage{parskip}
  }{% else
    \setlength{\parindent}{0pt}
    \setlength{\parskip}{6pt plus 2pt minus 1pt}}
}{% if KOMA class
  \KOMAoptions{parskip=half}}
\makeatother
\usepackage{longtable,booktabs,array}
\newcounter{none} % for unnumbered tables
\usepackage{calc} % for calculating minipage widths
% Correct order of tables after \paragraph or \subparagraph
\usepackage{etoolbox}
\makeatletter
\patchcmd\longtable{\par}{\if@noskipsec\mbox{}\fi\par}{}{}
\makeatother
% Allow footnotes in longtable head/foot
\IfFileExists{footnotehyper.sty}{\usepackage{footnotehyper}}{\usepackage{footnote}}
\makesavenoteenv{longtable}
\usepackage{graphicx}
\makeatletter
\newsavebox\pandoc@box
\newcommand*\pandocbounded[1]{% scales image to fit in text height/width
  \sbox\pandoc@box{#1}%
  \Gscale@div\@tempa{\textheight}{\dimexpr\ht\pandoc@box+\dp\pandoc@box\relax}%
  \Gscale@div\@tempb{\linewidth}{\wd\pandoc@box}%
  \ifdim\@tempb\p@<\@tempa\p@\let\@tempa\@tempb\fi% select the smaller of both
  \ifdim\@tempa\p@<\p@\scalebox{\@tempa}{\usebox\pandoc@box}%
  \else\usebox{\pandoc@box}%
  \fi%
}
% Set default figure placement to htbp
\def\fps@figure{htbp}
\makeatother
\setlength{\emergencystretch}{3em} % prevent overfull lines
\providecommand{\tightlist}{%
  \setlength{\itemsep}{0pt}\setlength{\parskip}{0pt}}
% Font setup
\usepackage{fontspec}
% TeX Gyre Termes: a Times-compatible serif under the GUST Font License,
% freely redistributable. The public repository must not embed subsets of a
% proprietary font in artifacts it distributes.
\setmainfont{TeX Gyre Termes}

% Page geometry and line spacing
\usepackage[margin=1in]{geometry}
\linespread{1.2}

% Color and section formatting
\usepackage{xcolor}
\definecolor{sectioncolor}{RGB}{140,45,15} % dark reddish-brown

% Colored links setup
\usepackage{hyperref}
\hypersetup{
  colorlinks=true,
  linkcolor=sectioncolor,
  urlcolor=blue,
  citecolor=sectioncolor
}

% Section formatting
\usepackage{titlesec}
\titleformat{\section}
  {\color{sectioncolor}\normalfont\Large\bfseries\sffamily\MakeUppercase}
  {}{0pt}{}
\titlespacing{\section}{0pt}{*4}{*1}
\usepackage{etoolbox}
\preto{\section}{\vspace{1em}\textcolor{sectioncolor}{\hrule}\vspace{0.5em}}

\titleformat{\subsection}
  {\color{sectioncolor}\normalfont\large\bfseries\sffamily\MakeUppercase}
  {}{0pt}{}
\titlespacing{\subsection}{0pt}{*3}{*1}

\titleformat{\subsubsection}
  {\color{sectioncolor}\normalfont\normalsize\itshape\sffamily}
  {}{0pt}{}

\usepackage{fancyhdr}
\pagestyle{fancy}
\fancyhf{} % Clear all header and footer fields

% Header
\lhead{SPPH 604: Application of Advanced Epidemiological Methods}
\rhead{SPPH 604}

% Footer
\renewcommand{\footrulewidth}{0.4pt} % Horizontal line above footer
\fancyfoot[L]{University of British Columbia}
\fancyfoot[R]{\thepage}
\usepackage{booktabs}
\usepackage{longtable}
\usepackage{array}
\usepackage{multirow}
\usepackage{wrapfig}
\usepackage{float}
\usepackage{colortbl}
\usepackage{pdflscape}
\usepackage{tabu}
\usepackage{threeparttable}
\usepackage{threeparttablex}
\usepackage[normalem]{ulem}
\usepackage{makecell}
\usepackage{xcolor}
\usepackage{bookmark}
\IfFileExists{xurl.sty}{\usepackage{xurl}}{} % add URL line breaks if available
\urlstyle{same}
% fallback for those not using the hyperref driver hyperxmp:
\makeatletter
\@ifundefined{xmpquote}{\newcommand{\xmpquote}[1]{#1}}{}
\makeatother
\hypersetup{
  pdftitle={SPPH 604: Application of Advanced Epidemiological Methods},
  hidelinks,
  pdfcreator={LaTeX via pandoc}}

\title{SPPH 604: Application of Advanced Epidemiological Methods}
\author{}
\date{\vspace{-2.5em}}

\begin{document}
\maketitle

\section{Acknowledgement}\label{acknowledgement}

UBC's Point Grey Campus is located on the traditional, ancestral, and
unceded territory of the xwməθkwəy̓əm (Musqueam) people. The land it is
situated on has always been a place of learning for the Musqueam people,
who for millennia have passed on their culture, history, and traditions
from one generation to the next on this site.

\section{Course Information}\label{course-information}

\begin{table}[H]
\centering
\begin{tabular}[t]{>{\raggedright\arraybackslash}p{3cm}>{\raggedright\arraybackslash}p{9cm}>{\raggedright\arraybackslash}p{2cm}}
\toprule
Course Title & Course Code \& Schedule & Credit Value\\
\midrule
\cellcolor{gray!10}{Application of Advanced Epidemiological Methods} & \cellcolor{gray!10}{Class: SPPH 604 001 (Tue 9:30 am - 12 pm unless otherwise specified): Zoom \newline Lab: SPPH 604 T01 (Thu 10 am - 12 pm): Room SPPH 143 \textcolor{red}{Class starting from Sept 10th and lab from Sept 17th, 2026} \newline Lab is in-person and class is Zoom sessions, unless specified otherwise.} & \cellcolor{gray!10}{3}\\
\bottomrule
\end{tabular}
\end{table}

\subsection{Course Description}\label{course-description}

\emph{As published in the UBC Academic Calendar:} Statistical approaches
for confounding, missingness, and complex surveys, the development of
analysis plans, the analysis and interpretation of real-world
epidemiologic data, and the communication of findings.

This course is not eligible for Credit/D/Fail grading.

\subsection{Course Modality}\label{course-modality}

This course is delivered in mixed mode. The Tuesday class (SPPH 604 001)
is \textbf{online}, held synchronously by Zoom; the Thursday lab (SPPH
604 T01) is \textbf{in person} in SPPH 143, where the hands-on work and
peer collaboration take place. Week 1 has no Tuesday session: the course
opens on \textbf{Thursday, September 10, 10 am--12 pm, by Zoom}, and the
in-person labs (L1--L6) begin \textbf{Thursday, September 17}. Zoom
links are posted on Canvas. Scheduled contact time is \textbf{4.5 hours
per week in Weeks 2--7} (usually a 2.5-hour Tuesday class plus a 2-hour
Thursday lab). In Weeks 8--12 the Tuesday class runs alone, as formal
instruction gives way to instructor-supported project time; the M1, M2
and M3 presentation sessions use the full 9 am--12 pm block. There is no
scheduled session in Week 10 (Midterm Break, Nov 9--11). See the
schedule table for the week-by-week breakdown.

\subsection{Course Instructor Biographical
Statement}\label{course-instructor-biographical-statement}

\begin{quote}
I am an Associate Professor in Health Data Science at the UBC School of
Population and Public Health
(\href{https://spph.ubc.ca/person/ehsan-karim/}{SPPH}), a Scientist at
the Centre for Advancing Health Outcomes, and an associate member of the
Department of \href{https://www.stat.ubc.ca/}{Statistics} at UBC. I was
previously a Michael Smith Health Research BC Scholar. I received my PhD
in Statistics from the University of British Columbia and completed my
postgraduate training in the Department of Epidemiology, Biostatistics,
and Occupational Health at
\href{https://www.mcgill.ca/epi-biostat-occh/}{McGill University}. My
current research focuses on causal inference, real-world observational
data analyses, and the application of machine learning approaches in
epidemiologic studies. This work is supported by grants from NSERC,
CIHR, and MS Canada.
\end{quote}

\subsection{Instructor Teaching
Philosophy}\label{instructor-teaching-philosophy}

\begin{quote}
I teach the application of epidemiological methods with a focus on
bridging theory and practice. My approach is student-centered, grounded
in active engagement, reproducibility, and real-world application. I
emphasize hands-on learning using real data, supporting students to
critically evaluate evidence and design rigorous and reproducible
analyses. Through open-access resources and detailed feedback, I aim to
foster technical mastery, independence, and scholarly confidence. I am
committed to accessibility so all learners can succeed and contribute
meaningfully to public health research. My goal is to cultivate critical
thinkers who question assumptions, assess bias, and apply
epidemiological principles with clarity and rigor. I strive to make
epidemiological thinking accessible and relevant across disciplines,
preparing students to collaborate effectively in diverse research and
policy settings.
\end{quote}

\subsection{Course Philosophy: The Capstone
Approach}\label{course-philosophy-the-capstone-approach}

\textbf{Scope \& The ``Need-to-Know'' Approach}: This course is designed
as a capstone experience, meaning its primary goal is to integrate and
apply your existing skills to a complete research project. It is not
intended to be an exhaustive theoretical survey of every aspect of
confounding, missing data, or complex survey analysis. Instead, we focus
on the application of these methods to real-world data.

\textbf{The ``Replicate -\textgreater{} Interrogate -\textgreater{}
Improve -\textgreater{} Defend'' Spine}: The core structure of this
course is built around a single, term-long, reproducible reanalysis
project based on a recent open-data paper (\emph{recent} is defined
under \textbf{Choosing Your Project (M0 Feasibility Gate)}, below). You
will replicate the authors' findings, interrogate their methodological
choices, propose and execute an improvement plan, and defend your
reasoning live.

We recognize that students enter this course with diverse backgrounds;
some may have gaps in computing (R), while others may have gaps in
epidemiological theory. This is expected. Our approach is
``just-in-time'' learning: rather than trying to master every nuance of
a topic in the abstract, we will address specific knowledge gaps as they
arise. We focus strictly on the tools and concepts you need to
successfully execute your final project and to competently peer-review
the work of others. The goal is not theoretical perfection, but
practical, reproducible competence in both your own analysis and your
critical evaluation of your colleagues' research.

\section{Prerequisites}\label{prerequisites}

All of:

\begin{itemize}
\tightlist
\item
  SPPH 500 Analytical Methods in Epidemiological Research
\item
  SPPH 502 Epidemiological Methods I
\item
  SPPH 503 Epidemiological Methods II
\end{itemize}

and one of:

\begin{itemize}
\tightlist
\item
  SPPH 400 Statistics for Health Research, \textbf{or}
\item
  SPPH 521 Statistics for Health Research
\end{itemize}

There are no co-requisites. Where this list and the UBC Academic
Calendar differ, the Calendar governs.

This is an applied, data-intensive course designed to bridge the gap
between epidemiological theory and real-world practice. Unlike
introductory courses, we prioritize ``just-in-time'' learning,
addressing specific computing or theoretical gaps as they arise during
your analysis. While we strongly encourage prior familiarity with R, we
adopt a ``Zero to AI Co-Pilot'' pedagogy. You are permitted to use
AI-augmented workflows to assist with data wrangling and programming. AI
should be a supplementary resource for learning; students remain fully
responsible for the integrity, auditability, and accuracy of all code
and statistical outputs.

\section{Contacts}\label{contacts}

\begin{table}[H]
\centering
\begin{tabular}[t]{>{\raggedright\arraybackslash}p{4cm}>{\raggedright\arraybackslash}p{5cm}>{\raggedright\arraybackslash}p{5cm}}
\toprule
Course Instructor(s) & Contact Details & Office Hours\\
\midrule
\cellcolor{gray!10}{\href{http://www.ehsank.com/}{Dr. M. Ehsan Karim}} & \cellcolor{gray!10}{ehsan dot karim at ubc dot ca \newline The typical response time is within 48 hours on weekdays.} & \cellcolor{gray!10}{By appointments, or as posted on canvas.}\\
\bottomrule
\end{tabular}
\end{table}

\section{Other Instructional Staff}\label{other-instructional-staff}

The teaching team includes GTAs: (TBA). To communicate with the
computing TA, come to weekly mandatory labs and specified office hours
(as announced in Canvas). Office hours are a time to ask for
clarification, discuss ideas, or receive guidance on assignments. If you
need further support, contact with the course instructor, and additional
office hours will be assigned accordingly. Do not email the TAs
directly. For any issues communicate with the instructor directly.

\section{Course Structure}\label{course-structure}

\subsection{Course Purpose}\label{course-purpose}

The purpose of this course is to provide students with learning
opportunities to implement fundamental epidemiological concepts through
the application of methods using population and public health datasets.
Students will also be exposed to emerging epidemiological methodologies
that are frequently being applied to population and public
health-related research questions.

\subsection{Workload \& Time Management}\label{workload-time-management}

This 3-credit graduate-level course provides \textbf{40.5 hours of
scheduled contact time} over the term -- see the schedule table for the
week-by-week breakdown -- exceeding the 39 hours conventionally
associated with a 3-credit Winter Session course. The syllabus itself
follows
\href{https://scs-senate-2021.sites.olt.ubc.ca/files/20260218_V-130_Approved-1.pdf}{UBC
Senate Policy V-130.1} (\emph{Content and Distribution of Course
Syllabi}, effective 01 September 2026). To prevent end-of-term workload
cliffs, the project work is spread across the term as five weekly
increments (P1--P5), each building on that week's lab rather than
waiting for the end of the course. The schedule is structured to evolve:
as you progress through the term, formal instruction gives way to
dedicated, instructor-supported project time.

\begin{itemize}
\tightlist
\item
  \textbf{Baseline for Success:} We anticipate that the scheduled class
  and lab time is sufficient for students to successfully meet learning
  objectives and achieve a strong standing.
\item
  \textbf{Prerequisites:} This estimate assumes mastery of the
  prerequisite materials. Students who are using this time to learn data
  analysis or review epidemiological theory from scratch should
  anticipate needing additional time to bridge those gaps.
\item
  \textbf{Exceptional Performance:} As per the course grading policy, an
  A+ is reserved for work that \emph{significantly exceeds}
  expectations. This is distinguished by depth of insight, creativity,
  and rigor, not merely by the number of hours logged. While some
  students may choose to invest extra time to refine their work to this
  level, the standard workload is designed to ensure you can produce
  high-quality, reproducible research without exceeding university time
  guidelines.
\item
  \textbf{Workflow:} Some weeks may require additional time depending on
  project milestones. Most major work is group-based -- M3 is assessed
  individually, within your group's improved-analysis presentation -- so
  students are expected to coordinate proactively to plan and divide
  tasks effectively.
\end{itemize}

\subsection{Choosing Your Project (M0 Feasibility
Gate)}\label{choosing-your-project-m0-feasibility-gate}

Your entire project is built on a single \textbf{recent, open-data
paper} that your group selects and ``locks'' at \textbf{M0 (Week 3)}.
\textbf{Recent} means published in the current calendar year or the five
preceding ones --- for 2026W, a publication year of \textbf{2021 or
later}; the demonstration paper, published in 2026, sits well inside
that window. Because the weekly increments deliberately exercise
specific methods --- cohort reconstruction, confounding, interaction,
design-aware estimation, and missing data --- not every paper can
support the full sequence. A paper is \textbf{eligible only if you can
identify all six} of the following; the M0 deliverable is a short
feasibility memo confirming each with a specific pointer (variable
names, file names, sample sizes):

\begin{enumerate}
\def\labelenumi{\arabic{enumi}.}
\tightlist
\item
  \textbf{Reconstructable data.} Accessible raw or near-raw microdata
  (e.g., NHANES, NHIS, or a comparable fully public source) from which
  you can rebuild the analytic cohort. A paper whose numbers can only be
  read from a summary table --- with no underlying data you can re-run
  --- does \textbf{not} qualify.
\item
  \textbf{A clear exposure and outcome.} A single primary association
  you can state in one sentence.
\item
  \textbf{Plausible confounders.} Measured covariates that support a
  defensible, causally-reasoned adjustment set.
\item
  \textbf{Adequate size and events.} Enough observations --- and, for
  time-to-event outcomes, enough events --- for the main analysis
  \emph{and} at least one subgroup/interaction analysis to be stable.
\item
  \textbf{Design variables.} The paper must contain a meaningful design
  or outcome-structure feature that this course teaches you to handle,
  and your memo must name it and point to the variables that encode it.
  Normally that is complex-survey sampling --- the weights, strata, and
  PSU identifiers that make a design-aware estimate possible. The other
  qualifying case is a time-to-event outcome, where the feature is the
  follow-up time and the censoring indicator and the paper's handling of
  them leaves something you can genuinely interrogate. Complex-survey
  data is the default case and the one the demonstration paper uses, but
  it is not the only one that qualifies. A paper whose data carry
  neither does not qualify, because P4 asks you to re-estimate the
  headline result the way the design demands and there would be nothing
  to re-estimate.
\item
  \textbf{A genuine incomplete-observation problem.} Point to at least
  one variable your locked analysis actually uses --- the exposure, the
  outcome, or a covariate in the main model --- that is absent for
  records which should have carried it. That is the problem P5 analyses,
  and it is what this criterion asks for. Your cohort will also shrink
  through designed subsampling, eligibility restrictions, the disease
  definition, and failure to link or follow up the outcome. You must be
  able to tell those apart from missingness, and at P5 you report them
  and scope your conclusion around them --- but none of them satisfies
  this criterion on its own, because P5 does not address them with the
  methods this course teaches. Do not argue from magnitude: a single
  covariate missing for a small share of the cohort clears this gate.
  The test is whether you could have known in advance that handling it
  would change nothing; if you could not, it is worth analysing. What
  fails is a paper whose analysis variables are complete by
  construction, leaving nothing to measure.
\end{enumerate}

\textbf{Criterion 5 is not waivable:} complex-survey sampling is the
normal route through it, a time-to-event structure whose handling you
can genuinely interrogate is the other, and a paper carrying neither has
nothing for P4 to re-estimate. Design features this course does not
teach you to handle --- clustering outside a survey design, repeated
measures --- do not clear it.

\textbf{Designed subsampling does not satisfy criterion 6.} Where a
study deliberately measured something in a subsample and supplies a
weight for it, that is a design feature with a documented fix --- the
right weight, plus an honest sentence naming who your estimate is about.
It is a scope question, not a missing-data problem, and it does not
clear this criterion.

Relatedly, the increment you build on criterion 6 must analyse a
missingness problem that \textbf{arises from your locked scientific
analysis} --- the absent value on a variable your locked model uses that
criterion 6 required you to point at. Introducing extra variables solely
to manufacture one is not a missing-data analysis; it is a change of
adjustment set wearing a disguise. Selection, eligibility and
outcome-linkage losses, designed subsampling included, are reported at
P5 and scoped around; they are not what this increment analyses.

The demonstration paper used throughout the course (a MASLD / NHANES
all-cause mortality analysis) passes all six and is used to illustrate
what each increment should look like. If your first choice fails the
gate, the teaching team will help you identify an eligible replacement
\textbf{before} the M0 lock; a paper that cannot clear all six criteria
cannot be carried into P1--P5.

\textbf{What the lock means.} Your paper locks at the M0 deadline shown
in the schedule table below: after that deadline there are no elective
changes of paper, and you carry the paper you named through P1--P5, M1,
M2, M3, and M4. The lock is not a commitment to an unworkable paper. If
the memo itself shows that the paper cannot clear all six criteria, the
teaching team may require you to replace it; the paper you then move to
becomes your locked paper. That is the teaching team's decision to make,
not a swap you can elect.

\section{Schedule of Topics}\label{schedule-of-topics}

Weekly materials will be posted in the Canvas ``Modules'' section by 9
AM each Monday, if any. Lab assignments are usually due by the end of
the lab. If you encounter technical issues during lab, please talk to
the TA immediately. Extensions will be considered on a case-by-case
basis.

\begingroup\fontsize{8}{10}\selectfont

\begin{longtable}[t]{>{\raggedright\arraybackslash}p{1.5cm}>{\raggedright\arraybackslash}p{1cm}>{\raggedright\arraybackslash}p{4cm}>{\raggedright\arraybackslash}p{4.5cm}>{\raggedright\arraybackslash}p{4cm}}
\toprule
Week \& Dates & Contact hrs (cum.) & Methods / In-class & Lab (fixed data) $\rightarrow$ Increment & Deadlines (date \& time)\\
\midrule
\cellcolor{gray!10}{\# 1 (Sept 10)} & \cellcolor{gray!10}{2.0} & \cellcolor{gray!10}{Course + project overview; R+GitHub setup} & \cellcolor{gray!10}{-} & \cellcolor{gray!10}{}\\
\# 2 (Sept 15 / 17) & 6.5 & Survey data sources; reading Methods critically & \textbf{L1} Data wrangling $\rightarrow$ \emph{shortlist papers} & \textbf{L1} due end of lab, Thu Sept 17\\
\cellcolor{gray!10}{\# 3 (Sept 22 / 24)} & \cellcolor{gray!10}{11.0} & \cellcolor{gray!10}{Analytic cohort: case definitions, eligibility \& the funnel; feasibility \& peer-review principles} & \cellcolor{gray!10}{\textbf{L2} Analytic data $\rightarrow$ \emph{P1: eligibility}} & \cellcolor{gray!10}{\textbf{M0}: Paper locked, due by end of Tue class, Sept 22 (12 PM) \newline \textbf{L2} due end of lab, Thu Sept 24}\\
\# 4 (Sept 29 / Oct 1) & 15.5 & Variable roles; confounding \& DAGs & \textbf{L3} Confounding $\rightarrow$ \emph{P2: Table 1} & \textbf{L3} due end of lab, Thu Oct 1\\
\cellcolor{gray!10}{\# 5 (Oct 6 / 8)} & \cellcolor{gray!10}{20.0} & \cellcolor{gray!10}{Interaction \& effect modification} & \cellcolor{gray!10}{\textbf{L4} Interaction $\rightarrow$ \emph{P3: effect modification}} & \cellcolor{gray!10}{\textbf{L4} due end of lab, Thu Oct 8}\\
\addlinespace
\# 6 (Oct 13 / 15) & 24.5 & Complex survey analysis & \textbf{L5} Survey analysis $\rightarrow$ \emph{P4: main estimate} & \textbf{L5} due end of lab, Thu Oct 15\\
\cellcolor{gray!10}{\# 7 (Oct 20 / 22)} & \cellcolor{gray!10}{29.0} & \cellcolor{gray!10}{Missing-data analysis} & \cellcolor{gray!10}{\textbf{L6} Missing data $\rightarrow$ \emph{P5: missingness}} & \cellcolor{gray!10}{\textbf{L6} due end of lab, Thu Oct 22}\\
\# 8 (Oct 27 / 29) & 32.0 & \textbf{M1 Replication presentation (Starts 9am)} & - & \textbf{M1}: Replication presented in class \newline slides + code zip due Tue Oct 27, 4 PM\\
\cellcolor{gray!10}{\# 9 (Nov 3 / 5)} & \cellcolor{gray!10}{35.0} & \cellcolor{gray!10}{\textbf{M2 Critique Session (Starts 9am)}: reviewer groups present, Q/A} & \cellcolor{gray!10}{-} & \cellcolor{gray!10}{\textbf{M2}: critique presented in class \newline slides due Tue Nov 3, 4 PM}\\
\# 10 (Nov 10 / 12) & 35.0 & (Midterm Break: Nov 9 - 11) & - & \\
\addlinespace
\cellcolor{gray!10}{\# 11 (Nov 17 / 19)} & \cellcolor{gray!10}{38.0} & \cellcolor{gray!10}{\textbf{M3 Improved-analysis presentation + individual oral defense (Starts 9am)}} & \cellcolor{gray!10}{-} & \cellcolor{gray!10}{\textbf{M3}: Live Defense, Tue Nov 17 (from 9 AM)}\\
\# 12 (Nov 24 / 26) & 40.5 & Open discussion & - & \textbf{Grade Re-weighting Plan} due Thu Nov 26, 4 PM\\
\cellcolor{gray!10}{Final deadline (Fri Dec 11)} & \cellcolor{gray!10}{-} & \cellcolor{gray!10}{} & \cellcolor{gray!10}{-} & \cellcolor{gray!10}{\textbf{M4}: Final letter + code zip due Fri Dec 11, 4 PM \newline \textcolor{red}{\textbf{Bonus}}: journal-submission pledge, same deadline}\\
\bottomrule
\end{longtable}
\endgroup{}

\emph{Note: This is a \textbf{proposed} schedule of topics; topic order
may be adjusted, but the deadlines above are fixed. Week 1 has no
Tuesday session -- the course opens Thursday, September 10, 10 am--12
pm, by Zoom. The in-person Thursday lab runs in Weeks 2--7 (L1--L6),
beginning September 17; the L0 setup materials are self-study and are
posted on Canvas. In Weeks 8--12 there is no scheduled Thursday session;
the Tuesday class absorbs the presentation milestones and the remaining
project work is instructor-supported time arranged with the teaching
team. The ``Contact hrs (cum.)'' column is the running total of
scheduled contact hours to the end of that week.}

\section{Learning Outcomes}\label{learning-outcomes}

By the end of this course, students will be able to:

\subsection{Study Design and Protocol
Development}\label{study-design-and-protocol-development}

\begin{itemize}
\tightlist
\item
  Critically appraise and refine epidemiologic research questions to
  ensure specificity and feasibility.
\item
  Develop structured study protocols aligned with public health
  objectives, incorporating principles of transparency and
  reproducibility.
\item
  Identify and justify appropriate study designs for varied research
  contexts.
\end{itemize}

\subsection{Data Analysis and
Interpretation}\label{data-analysis-and-interpretation}

\begin{itemize}
\tightlist
\item
  Translate epidemiologic questions into a reproducible statistical
  analysis plan.
\item
  Generate, interpret, and critique key statistical outputs, including
  model coefficients, confidence intervals, and diagnostics.
\item
  Evaluate the robustness of findings through sensitivity analyses and
  appropriate methodological critique.
\end{itemize}

\subsection{Communication and Scientific
Reasoning}\label{communication-and-scientific-reasoning}

\begin{itemize}
\tightlist
\item
  Communicate scientific reasoning and findings in clear, structured
  written formats suitable for academic and applied audiences.
\item
  Provide and integrate constructive peer feedback to improve clarity
  and scientific quality in collaborative work.
\item
  Synthesize technical and contextual elements of epidemiologic research
  in a capstone deliverable.
\end{itemize}

\subsection{How these outcomes are
met}\label{how-these-outcomes-are-met}

The section that follows maps the activities to these outcome clusters;
three of the outcomes deserve a more precise note. \emph{Structured
study protocols}: the M0 feasibility memo is where you must name in
advance, with specific pointers, the question, the exposure and outcome
definitions, the confounders, the design features and the missingness
your analysis will have to handle; the reproducible repository, together
with the M3 requirement that any later change to those decisions be
identified and justified, is what supplies the transparency and
reproducibility; no separate study-protocol document is written or
graded. \emph{Identifying and justifying study designs}: every project
reanalyses public microdata from a study someone else designed, so you
reconstruct and defend a design -- eligibility rules, comparison groups,
and the design features the data carry -- rather than choose one for a
new question; beyond your own locked paper, the contexts you meet are
the fixed teaching datasets of L1--L6, a national survey and a clinical
cohort, together with the second paper you appraise at M2.
\emph{Translating a question into a reproducible analysis plan}: the
plan is not a separately submitted document; the M0 memo states it in
outline, and your repository is its executable form, which is why your
code must run start to finish on a fresh machine and produce every
number you report.

\section{Learning Activities}\label{learning-activities}

The activities below are the means by which you achieve and demonstrate
the learning outcomes listed above. In broad terms: \textbf{M0 and
L1--L6 with P1--P5} build \emph{Study Design and Protocol Development}
and \emph{Data Analysis and Interpretation}; \textbf{M1 and M3} assess
\emph{Data Analysis and Interpretation}, including robustness and
methodological critique; \textbf{M2 and M4} build and assess
\emph{Communication and Scientific Reasoning}, with M2 targeting the
outcome on giving and integrating constructive peer feedback and M4
serving as the capstone deliverable. Sustained, critical reading of your
locked paper, its supplementary material, and the relevant chapters of
the two open-access textbooks is assumed throughout and is what makes M0
and M2 possible; specific chapters are signposted week by week in
Canvas.

\begin{itemize}
\tightlist
\item
  \textbf{Weekly lab data analysis activities (L1-L6, P1-P5)}: Students
  will complete in-lab data analysis exercises (L1-L6) applying concepts
  on fixed datasets; the matching project increments (P1-P5) apply those
  same concepts to their own chosen paper. The project increments
  \textbf{P1--P5 are optional and ungraded}: they carry no independent
  weight, no separate deadline, and nothing is submitted for them. They
  are the weekly building blocks of your replication and are assessed
  only once they are compiled and submitted as \textbf{M1} --- and the
  practice they give you is what you draw on to critique another group's
  M1 at M2 and to defend your own analysis at M3.
\item
  \textbf{M1 \& M2 (Slide Decks; code zip with M1 only)}: To combat
  AI-generated essay bloat, M1 and M2 are submitted strictly as
  presentation slide decks. \textbf{M1} is additionally accompanied by a
  \textbf{zip archive of your reproducible repository, downloaded from
  GitHub} (you are not required to grant us repository access);
  \textbf{M2 is slides only}, because you do not write code for the
  other group's paper. Traditional written reports are not accepted for
  these milestones. \textbf{M1} is your five project increments compiled
  into one replication-and-interrogation deck, closing with your group's
  judgment of the \textbf{single threat that most endangers the paper's
  conclusion}; it is presented and submitted on \textbf{Tuesday, October
  27, due 4 PM}, and the deck and code zip are released to its assigned
  reviewers as soon as they are in.
\item
  \textbf{M2 (Critique of another group's M1)}: In the Tuesday Week 9
  class your group presents a critique of the M1 you were assigned, and
  submits the slides by \textbf{4 PM the same day}. Every member
  presents a section and answers questions on it. Pairings are announced
  when M1 is submitted, and run in a cycle so that no group reviews the
  group reviewing it. The scope is fixed at three things, in about six
  slides: \textbf{(1) did their code run}, and if not, how far you got
  and where it stopped; \textbf{(2) two or three substantive points
  about the epidemiology} --- the exposure and outcome definitions, the
  adjustment set and how it was justified, selection into the cohort,
  whether the stated interpretation follows from the estimate they
  actually report, and how far the finding would generalize; and
  \textbf{(3) whether you agree with their choice of dominant threat}
  --- and if not, which threat you think dominates and why. You are not
  expected to have arrived knowing the clinical literature of a paper
  you did not choose, but you are expected to go and read enough of it
  to make the point you are making --- the paper's citations and the
  code in front of you are within reach. Put a point as a question only
  after you have looked and the answer is genuinely unsettled, and then
  say what you checked: ``the paper cites no source for this threshold''
  is a finding; ``we are not clinicians'' is not. A critique that finds
  little wrong is entirely acceptable when it is evidenced: a reviewer
  who recommends acceptance with reasons is doing the job. You are
  graded on the \textbf{quality of the critique} --- whether you
  identify real methodological weaknesses, express them professionally,
  and suggest a practical way forward --- never on how many faults you
  produce. You do \textbf{not} write code for the other group's paper
  and you do \textbf{not} implement anything on it; your recommendations
  are input that they must answer at M3.
\item
  \textbf{M3 (Live Oral Defense)}: M3 is where you \textbf{improve and
  defend your analysis}: address the most important points raised at M2,
  identify any further changes your group made independently, and
  justify the final analytic decisions. Your M1 analysis is the starting
  point, and M2 is an input to it --- not its only source of insight.
  You are \textbf{not} obliged to implement your reviewers'
  recommendations; you are obliged to answer them, and you are free to
  change things they never raised. Any change to the analytic
  population, estimand, adjustment set, study design, variable
  definitions, missing-data strategy or model specification must be
  \textbf{identified and justified}, whether a reviewer prompted it or
  you did. The presentation shows \emph{what} you did; the defense is
  about \textbf{whether it was justified}, and whether the change alters
  (or fails to alter) the interpretation in a way you can explain and
  stand behind. Grading rewards sound epidemiological reasoning under
  questioning, not whether your code merely runs.
\item
  \textbf{M4 (Final Reanalysis Letter + code zip)}: Your final project
  will be a short, densely written reanalysis letter/commentary
  (\textasciitilde1,000--1,500 words) accompanied by a \textbf{zip
  archive of your fully reproducible repository, downloaded from
  GitHub}.
\item
  \textbf{Reproducibility}: In health data science, code that only runs
  on your machine is considered unfinished. We grade strictly on
  auditability. Your code must run from start to finish on a fresh
  machine without manual intervention.
\end{itemize}

\section{Learning Materials}\label{learning-materials}

\textbf{Estimated cost of required learning materials and activities:
\$0.} Both required textbooks are open access, the required
\texttt{svyTable1} package is free, R / RStudio / R Markdown are free of
charge, and a GitHub account is free. There are no lab fees, field-trip
costs, course packs, or paid online assessment tools. Students supply
their own laptop (see \emph{Learning Resources}); if access to a
suitable computer is a barrier for you, contact the instructor.

We have 2 open access textbooks.

\begin{itemize}
\tightlist
\item
  Karim M. E., Epi-OER team (2024). Advanced Epidemiological Methods.
  \url{https://ehsanx.github.io/EpiMethods/}.
\item
  Karim ME, Jeong D, Yusuf F (2021) Scientific Writing for Health
  Research.
  \url{https://ehsanx.github.io/Scientific-Writing-for-Health-Research/}.
\end{itemize}

Open software packages:

\begin{itemize}
\tightlist
\item
  Karim E (2025). svyTable1: Create Survey-Weighted Descriptive
  Statistics and Diagnostic Tables. R package version 0.14.1,
  \url{https://ehsanx.github.io/svyTable1}.
\end{itemize}

The following textbooks are suggested for further reading and additional
learning (also available via UBC library; \emph{no need to buy}):

\begin{itemize}
\tightlist
\item
  Heard, S. B. (2016). The scientist's guide to writing: How to write
  more easily and effectively throughout your scientific career.
  Princeton University Press.
\item
  Heeringa, S.G., West, B.T., Berglund, P.A (2010) Applied Survey Data
  Analysis, Taylor and Francis, Florida.
\end{itemize}

\section{Assessments of Learning}\label{assessments-of-learning}

The course will involve lectures and course assignments focused on the
application of epidemiological methods. Although students will gain some
expertise in statistical computation and programming, this course is
focused on the application of epidemiologic analytic methods. Students
will be evaluated based on the following elements:

\begin{enumerate}
\def\labelenumi{(\alph{enumi})}
\tightlist
\item
  apply key epidemiologic concepts,
\item
  understanding of analytic approaches to reduce study biases,
\item
  the application of epidemiological methods to population and public
  health research questions, and
\item
  the appropriate interpretation of analytic estimates from analytic
  output.
\end{enumerate}

\textbf{Students must submit a ``Grade Re-weighting Plan'' via Canvas
inbox message to the Instructor by Thu Nov 26, 4 PM (Week 12),
confirming which assignments are being shifted to the Final Project (M4)
or to M1, and stating a brief reason or attaching an approved
accommodation. That is the single deadline for all weight-shifting
requests.} Any Bonus points earned will be added to this total, up to a
maximum final grade of 100\%.

\begin{table}[H]
\centering
\begin{tabular}[t]{>{\raggedright\arraybackslash}p{10cm}>{\centering\arraybackslash}p{2cm}}
\toprule
Assessment & Weight\\
\midrule
\cellcolor{gray!10}{Labs L1–L6, Complete/Incomplete (Shiftable to M1)} & \cellcolor{gray!10}{10\%}\\
P1–P5 increments (optional, ungraded; compiled into M1) & 0\%\\
\cellcolor{gray!10}{M0 feasibility gate} & \cellcolor{gray!10}{0\%}\\
M1 replication + interrogation: P1–P5 compiled + dominant-threat judgment (Shiftable to M4) & 30\%\\
\cellcolor{gray!10}{M2 Critique of another group's M1 (not shiftable)} & \cellcolor{gray!10}{15\%}\\
\addlinespace
M3 defense: improved analysis + response to review (Shiftable to M4) & 15\%\\
\cellcolor{gray!10}{M4 reanalysis letter + reproducible repo} & \cellcolor{gray!10}{30\%}\\
\textcolor{red}{Bonus} [Journal-submission pledge] & \textcolor{red}{+2\%}\\
\bottomrule
\end{tabular}
\end{table}

\emph{Note: Shifting components to M4 requires a valid reason or formal
accommodation request.} \emph{Note: The weekly project increments
(P1--P5) are \textbf{optional and ungraded} --- they carry no
independent weight, no separate deadline, and are not separately
submitted. They are the building blocks of M1, into which they are
compiled. M0 is a pass/fail feasibility gate (0\%).}

Peer-review quality at M2 will be judged by whether the reviewer can
identify potential weaknesses in the work they are reviewing, express
them professionally, and suggest a practical way forward (if possible).
A review that finds little to correct is judged on the evidence it gives
for that conclusion, not penalized for the absence of faults.
Assessments are designed to build incrementally toward the final
project. Peer reviews provide structured feedback loops to refine your
analysis and writing. The bonus is earned by submitting, alongside M4, a
signed pledge committing to submit the work to a peer-reviewed journal
by January 31, 2027. This pledge is due with M4 (Dec 11, 4 PM) and the
deadline is non-negotiable.

\subsection{Assignment Code Table}\label{assignment-code-table}

\begin{table}[H]
\centering
\begin{tabular}[t]{ll}
\toprule
Assignment.Code & Description\\
\midrule
\cellcolor{gray!10}{M0} & \cellcolor{gray!10}{Paper locked}\\
L1-L6 & Lab data analysis\\
\cellcolor{gray!10}{P1-P5} & \cellcolor{gray!10}{Project increments}\\
M1 & P1-P5 compiled + threat judgment\\
\cellcolor{gray!10}{\textcolor{blue}{M2}} & \cellcolor{gray!10}{Critique of another group's M1}\\
M3 & Live defense\\
\cellcolor{gray!10}{\textcolor{blue}{M4}} & \cellcolor{gray!10}{Final letter + repo}\\
\bottomrule
\end{tabular}
\end{table}

\emph{Note: Permitted generative-AI use is specified per assessment in
the \textbf{Generative AI (GenAI) in Student-Submitted Work} section.
M0, M1, M2 and M4 each require a cover note containing a per-member
contribution note, any external-feedback disclosure, and an AI Use
Statement; it goes in the submitted document, and at M1 and M4 also as
\texttt{COVER-NOTE.md} in the root of the repository zip.}

\subsection{Grading}\label{grading}

The weekly labs (L1--L6) are graded \textbf{Complete / Incomplete}. A
lab is Complete when your knitted document ran without errors, carries
your group name and the first names of the members present, and was
submitted before the lab session ended. There is no band scale and no
rubric grid for labs: the point of a lab is that you leave the room able
to run the workflow.

The lab component is worth 10\% of the final grade and is computed from
your \textbf{best five of the six labs}. Each Complete among those best
five is worth \textbf{2 percentage points}, and there is no partial
credit within a lab. Six Completes and five Completes therefore both
earn the full 10\%; four Completes earn 8\%; three earn 6\%; two earn
4\%; one earns 2\%; and none earns 0\%. Lab credit is individual rather
than group-level: a lab is Complete for the members present in the room
and named on the submission (see \textbf{Late Assignments}, below), so
the best five are counted student by student.

Dropping your weakest lab is built into the scale, not something you
apply for. \textbf{One Incomplete costs you nothing}, whatever the
reason, and needs no request, no documentation and no \textbf{Grade
Re-weighting Plan}. The buffer covers exactly one lab and no more. If
you miss a second or subsequent lab, the labs carry the \emph{(Shiftable
to M1)} tag in the assessment table above, so you may ask to move the
weight of those further labs to M1 under the \textbf{Flexible Weighting
Policy} below, on the single \textbf{Grade Re-weighting Plan} deadline
of \textbf{Thu Nov 26, 4 PM} and subject to the approval that policy
requires. Accommodations arranged through the Centre for Accessibility
and academic concessions under the Academic Calendar are unaffected by
the buffer and remain available to you.

Milestones M1--M4 are graded on a coarse \textbf{3-band scale (Weak /
Adequate / Exceptional)} designed to prioritize objective computational
reproduction and live epidemiological reasoning over generic essay
polish. Detailed grading rubric grids will be provided via Canvas for
each of those milestones. \textbf{M0} is a pass/fail feasibility gate,
marked against its own checklist rather than on this scale, and the
P1--P5 increments are ungraded.

\begin{itemize}
\tightlist
\item
  \textbf{Exceptional (90\% to 100\%):} Demonstrates deep insight,
  accurately traces discrepancies to their root cause, or engages
  masterfully with counter-arguments.
\item
  \textbf{Adequate (75\% to 89\%):} Correctly executes the mechanics of
  the task and demonstrates sound epidemiological understanding.
\item
  \textbf{Weak (0\% to 74\%):} Generic or superficial; not specific
  enough to your data, your paper and your estimand; asserted without
  support from output you actually produced; technically incorrect; or
  disconnected from the estimand and the assumptions it rests on. Any
  one of these is enough on its own.
\end{itemize}

\subsection{Grade Appeal Policy}\label{grade-appeal-policy}

If you believe a grade is inaccurate or an administrative error
occurred, you may submit a written appeal. Appeals must adhere to the
following strict guidelines:

\begin{itemize}
\tightlist
\item
  Eligibility: Each assignment may be appealed only once.
\item
  Format: Maximum 2 pages, Times New Roman 11 pt, standard margins,
  submitted as a .docx file.
\item
  Content: clearly state the claim, cite specific rubric items/page
  numbers, and provide evidence/calculations supporting the request.
\item
  Process: The instructor and TAs will review the claim together. Peer
  reviewers will not be consulted. The teaching team may assign a single
  joint grade to the group if necessary.
\item
  Deadline: Appeals must be submitted within one week of the grade
  release.
\end{itemize}

\subsection{Formative Assessment \& Feedback
Policy}\label{formative-assessment-feedback-policy}

The interim milestones (M1, M2) are designed to help you build your
Final Project incrementally.

\textbf{M1 (Week 8)}: Your M1 replication is presented in class on
\textbf{Tuesday, October 27}, and the slide deck and code zip are due
the same day at \textbf{4 PM}. Marking is completed that week, so you
have your M1 feedback in hand before you build M2, which is presented
the following Tuesday.

\textbf{M2 (Week 9)} follows the same present-and-submit pattern: your
critique slides are presented in the Tuesday (Week 9) class and
\textbf{submitted by 4 PM that day}. There is no separate written round
afterwards --- the slides you present are the deliverable.

\subsubsection{Peer Review}\label{peer-review}

Peer review (M2) is an integral part of this course. The ability to
critically evaluate others' work is an essential skill for PhD students.
Reviews are evaluated based on clarity, constructiveness, and the
ability to spot methodological gaps. The M2 critique is graded by the
teaching team, not by your peers, and the critique you receive carries
no score: it is feedback, and your answer to it is what M3 assesses.

\subsubsection{Late Assignments}\label{late-assignments}

Group assignments (L1-L6 and the M1/M2 presentation decks) and the final
letter + repository (M4) must be submitted via Canvas. M3 is assessed
live and has no separate submission. Late milestone work is handled
\textbf{per milestone}, because the five milestones differ in kind:

\begin{itemize}
\tightlist
\item
  \textbf{M0} carries no percentage weight and is marked pass/fail, so
  lateness attracts no percentage penalty: a memo arriving after the
  deadline in the schedule table is still read and still marked. What
  the deadline fixes is the paper. After it there are no elective
  changes -- you may not decide to move to a different paper -- and if
  the memo shows that the paper you proposed is not feasible, the
  teaching team may require a replacement; the replacement we direct
  then becomes your locked paper. A late memo is still costly, because
  everything after M0 -- the optional P1-P5 increments, and then M1 --
  assumes a settled paper, so tell the instructor as soon as you know
  your memo will be late.
\item
  \textbf{M1} and \textbf{M2} are presented live and submitted the same
  day at 4 PM, so missing the presentation and submitting the deck late
  are two different failures and either can happen without the other.
  \emph{Not presenting} is an absence from a scheduled session, and is
  handled under the Flexible Weighting Policy below -- for M2, under
  that policy's M2 exception, because a missed critique leaves another
  group without feedback it is owed before M3. \emph{Presenting on time
  but submitting the deck after 4 PM} is a submission problem: the deck
  is accepted up to 24 hours after the deadline with a deduction of 10\%
  of the possible grade, and is not accepted after that, absent academic
  concession or an approved accommodation. The cost is not only yours --
  a late M1 reaches its assigned reviewers late, and late M2 slides
  delay feedback the reviewed group needs before M3.
\item
  \textbf{M3} has no submission and therefore cannot be late. It can
  only be missed, which is an absence from a scheduled session and is
  handled under the Flexible Weighting Policy below.
\item
  \textbf{M4} is the one milestone with a sliding late penalty, and it
  is set out in the M4 late-penalty paragraph below.
\end{itemize}

Nothing in this section limits academic concession or an approved
accommodation.

If you encounter technical issues during lab (for submitting L1-L6),
notify the TA immediately; extensions may be considered on a
case-by-case basis. Grades for lab exercises (L1-L6) are awarded to
students present in the room for the Thursday lab. Grades for the
presentation components (M1, M2, M3) are awarded to students who take
part in the scheduled Tuesday session. If your connection fails during a
presentation or defense, notify the instructor immediately; the missed
portion is completed at an alternative time and is not treated as an
absence. We recognize that graduate students may occasionally face
unavoidable scheduling conflicts. The Flexible Weighting Policy exists
to ensure you are not penalized for such absences.

The final letter + repository (M4) submitted after the due date and time
will be penalized 10\% of the possible grade, with an additional 5\%
deduction for each calendar day it is late. Extensions may be granted in
extenuating circumstances with the instructor's approval and appropriate
documentation (e.g., a medical certificate for the affected group
member).

\subsubsection{Flexible Weighting (``Safety Net'')
Policy}\label{flexible-weighting-safety-net-policy}

We recognize that graduate students manage complex professional and
personal schedules. To accommodate this, students may apply to transfer
the weight of missed interim assignments to the Final Project (M4) or
M1, as designated in the assessment table. To keep the course on a fixed
timeline, the in-term option available from the instructor is
reallocation of the weight of a missed component (as described below)
rather than resubmissions, retakes, or make-up sessions. This is the
instructor's in-term mechanism only. Accommodations arranged through the
Centre for Accessibility (see \emph{Academic Accommodation Letter},
below) and academic concessions that lie beyond the instructor's
authority -- for example deferred standing or late withdrawal, which are
granted by your graduate program or the Dean under the Academic
Calendar's
\href{https://vancouver.calendar.ubc.ca/campus-wide-policies-and-regulations/academic-concession}{academic
concession} regulations -- are unaffected by this policy and remain
available to you.

\begin{itemize}
\tightlist
\item
  \textbf{Mechanism}: If you are unable to contribute to a specific
  assignment, you may omit your name from the submission. You will not
  receive a zero; instead, the percentage weight of that assignment will
  be reallocated.
\item
  \textbf{Eligibility}: This applies to group assignments and
  presentations with a ``(Shiftable)'' tag. The M2 critique is excluded
  from routine weight-shifting, because a missed critique deprives
  another group of feedback they are entitled to before M3. Where
  academic concession or an approved accommodation applies to M2, the
  instructor will either assign an alternative review task or reallocate
  the M2 weight, at the instructor's discretion -- contact the
  instructor as soon as you know.
\item
  \textbf{Conditions}: Weight shifting requires instructor approval via
  a formal accommodation request or valid reason. \textbf{Transferred
  weight may be transferred again.} Weight moved from the labs to M1 may
  subsequently travel with M1 to M4; it does not become fixed to M1 on
  arrival. The ceiling therefore follows from M2 alone, which never
  moves: at most \textbf{85\%} of the grade can rest on the final letter
  -- M4's own 30\%, plus M1's 30\%, M3's 15\%, and the labs' 10\% by way
  of M1. Pushing the full 85\% onto M4 is allowed, but it is an
  extremely high-risk gamble.
\end{itemize}

For group assignments, individuals who do not participate will not
receive the group grade; instead, the assessment may be folded into
their Final Paper grade. There is a single standard and a single
deadline for this: submit a \textbf{Grade Re-weighting Plan} by
\textbf{Thu Nov 26, 4 PM (Week 12)}, listing the assignment(s) affected
and a brief reason, and obtain the instructor's approval.
Non-participation will be determined based on self-identification, peer
evaluations, and/or the discretion of the teaching team (instructor/TA).
Where you become aware of the need only after that deadline, contact the
instructor immediately; nothing in this policy limits your rights under
the Academic Calendar's academic concession regulations.

\section{University Policies}\label{university-policies}

UBC provides resources to support student learning and to maintain
healthy lifestyles but recognizes that sometimes crises arise and so
there are additional resources to access including those for survivors
of sexual violence. UBC values respect for the person and ideas of all
members of the academic community. Harassment and discrimination are not
tolerated nor is suppression of academic freedom. UBC provides
appropriate accommodation for students with disabilities and for
religious and cultural observances. UBC values academic honesty and
students are expected to acknowledge the ideas generated by others and
to uphold the highest academic standards in all of their actions.
Details of the policies and how to access support are available on the
\href{https://senate.ubc.ca/policies-resources-support-student-success}{UBC
Senate} website.

\section{Other Course Policies}\label{other-course-policies}

\subsection{Communication}\label{communication}

Email communication with the instructor should be reserved for
administrative matters. Email responses can be expected within 48 hours
on weekdays. Please include `SPPH 604' in the subject line. It's
strongly discouraged to email for assistance with the \emph{content} of
an assignment within 48 hours of its deadline; instead, seek help as
early as possible. This does \textbf{not} apply to questions about
academic integrity, attribution, collaboration boundaries, or what
generative-AI use is permitted -- ask those at any time.

\subsection{Academic Accommodation
Letter}\label{academic-accommodation-letter}

If you have an Academic Accommodation Letter from UBC's
\href{https://students.ubc.ca/about-student-services/centre-for-accessibility}{Centre
for Accessibility}, please send it to the instructor within the first 2
weeks of the class start. If you are not yet registered with the Centre
for Accessibility, you can begin that process at the same link. If there
is a delay in obtaining your letter, please simply notify the instructor
before the two-week mark so we can anticipate your needs.

This course is committed to inclusive and respectful learning. Please
speak with the instructor directly if you encounter barriers to
participation or have suggestions for improving accessibility. The
instructor welcomes constructive informal feedback on course
accessibility, pace, or delivery style. If there are ways the instructor
can make the course more inclusive or supportive of your learning, feel
free to reach out.

\subsection{Academic Integrity}\label{academic-integrity}

\subsubsection{Why this matters in this
course}\label{why-this-matters-in-this-course}

SPPH 604 exists because published epidemiological findings are supposed
to be independently reproducible --- and often are not. For a whole term
you will hold someone else's published work to that standard: rebuilding
their cohort, re-running their analysis, and reporting honestly whether
the numbers come out the same. That exercise only means something if
your own work can survive the same scrutiny. An audit you fabricated, a
replication you did not actually run, or a result you cannot account for
is not a small procedural lapse in this course; it is the exact failure
the course was designed to teach you to detect. Academic integrity here
is therefore not separate from the methods content --- it \emph{is} the
methods content.

\subsubsection{What is expected of you}\label{what-is-expected-of-you}

\begin{itemize}
\tightlist
\item
  Every number, figure, and table you submit was produced by code in
  your repository, and that code runs start to finish on a fresh machine
  (see \textbf{Reproducibility}, above).
\item
  Everything in your submission that you did not write yourself ---
  code, text, figures, an analytic idea --- is attributed to its source,
  whether that source is a person, a published paper, a package, a
  textbook, or a generative AI tool.
\item
  Your name on a group submission is a claim that you contributed to it
  (see \textbf{Authorship and contribution}, below).
\item
  When you are unsure whether something is allowed, you ask
  \textbf{before} you submit. Asking is never penalized. See
  \textbf{Where to get help}, below.
\end{itemize}

\subsubsection{Collaboration: what is and is not
permitted}\label{collaboration-what-is-and-is-not-permitted}

This course is almost entirely group-based, your code lives in a shared
GitHub repository, you review another group's unpublished work at M2,
and you defend your own work individually at M3. Each of those creates a
specific boundary. Anything not listed below, ask about.

\textbf{Within your own group --- unrestricted.} Write code together,
edit each other's text, pair-program, divide and re-divide the work
however you like. Group members are co-authors of M1, M2 and M4.

\textbf{Between groups --- ideas yes, content no.} You may discuss
methods, argue about a DAG at the whiteboard, point another group toward
a package or a vignette, and tell them how you solved a problem. You may
\textbf{not} send, receive, copy, or paste another group's code, slides,
text, output, or repository contents, and you may not submit work
another group produced. If a conversation with another group materially
shaped your analytic decision, say so in your cover note. That is a
strength, not a confession.

\textbf{The paper you are replicating --- reuse it, but cite it.} Many
open-data papers post replication code on GitHub, OSF, or as
supplementary material. \textbf{You may read, run, and adapt it} ---
locating and understanding the authors' code is legitimate replication
work, and pretending you found the cohort definition yourself when you
lifted it from their script is not. The condition is attribution: cite
the source in your repository \texttt{README} and identify, in your
cover note, which scripts or functions are derived from the original
authors' code and what you changed. Submitting the authors' pipeline
unchanged as your own replication, without saying so, is misconduct.

\textbf{Course, package, and public code --- reuse it, and cite it.}
Code from \emph{Advanced Epidemiological Methods}, \emph{Scientific
Writing for Health Research}, the \texttt{svyTable1} package and its
vignettes, CRAN documentation, Stack Overflow, or any public repository
may be reused freely. Attribute it with a URL in a code comment or in
your \texttt{README}. Routine idiom (a \texttt{ggplot} call, a
\texttt{dplyr} pipeline) needs no citation; a borrowed function,
algorithm, or non-obvious block does.

\textbf{Prior cohorts --- not permitted.} Using a previous year's SPPH
604 repository, slide deck, lab solution, or reanalysis letter as a
source for your own work is an unauthorized aid, whether or not it is
publicly posted. What you may republish is a separate question, settled
in \textbf{Copyright}: the handouts and other materials in the public
course repository are CC BY 4.0 and you may repost them with
attribution, while Canvas-only material, the exercise datasets, session
recordings and other students' work are not yours to publish.

\textbf{Repository visibility.} Keep your project repository
\textbf{private for the duration of the term}. Once final grades are
posted you are welcome --- and encouraged --- to make it public; ask if
you want to release it earlier, for example alongside a preprint.

\textbf{People outside your group.} Governed by the External Feedback
Policy below.

\textbf{Generative AI.} Governed by \emph{Generative AI (GenAI) in
Student-Submitted Work}, below.

\subsubsection{Peer review (M2): confidentiality and
non-appropriation}\label{peer-review-m2-confidentiality-and-non-appropriation}

At M2 you receive another group's unpublished M1 replication and its
code repository, several weeks before your own M3 and M4 are due. That
is a deliberate design choice and it carries three obligations:

\begin{itemize}
\tightlist
\item
  \textbf{No:} another group's unpublished slides, code, output or
  review materials, in whole or in part, pasted, uploaded, screenshotted
  or retyped (see \emph{Peer review (M2): confidentiality and
  non-appropriation}, above, and the \textbf{M2} row of the table); and
  any personal, confidential, or restricted data, should you ever
  encounter it in adjacent work.
\item
  \textbf{Not yours to use.} You may not copy or adapt a reviewed
  group's code, text, figures, output, or distinctive implementation
  into your own M2, M3, or M4 --- crediting them does not make it
  permitted. A general methodological insight is different: if reviewing
  their work shows you something you should reconsider in your own
  analysis, you may act on it, provided you justify and implement the
  change yourself, on your own paper, and disclose the influence in your
  cover note.
\item
  \textbf{Yours to write.} Your review group writes it together --- M2
  is a group submission, and dividing the sections among you is
  expected. What you may not do is reach outside that group: do not
  delegate any part of the review to someone who is not in it, and do
  not agree with another reviewing group on what to find or how hard to
  press. Reviewers award no marks --- the teaching team grades your
  critique, and the critique you receive carries no score --- so there
  is nothing to trade. Write the review you actually believe.
\end{itemize}

\subsubsection{Authorship and
contribution}\label{authorship-and-contribution}

Attaching your name to a group submission is a claim that you
contributed to it, and that claim is consequential: participation
determines whether you receive the group grade (see \textbf{Flexible
Weighting}). Leaving your name on work you did not contribute to
misrepresents your contribution and is an integrity matter, not merely a
workload one --- the Flexible Weighting Policy exists precisely so that
you never have to make that choice.

Accordingly, each group submission (M0, M1, M2, M4) includes a
\textbf{one-line contribution note per member} in the cover note --- for
example, \emph{``AB: cohort reconstruction, survey weighting code. CD:
DAG, confounder set, Table 1. EF: slides, missing-data sensitivity
analysis.''} This takes two minutes, gives the M3 individual defense a
documented basis, and is exactly the contributorship statement you will
be asked for if you act on the M4 journal-submission pledge. Peer
evaluations of participation are also to be given honestly; they inform,
but do not solely determine, participation decisions.

\subsubsection{What counts as misconduct in this
course}\label{what-counts-as-misconduct-in-this-course}

Most of you have been trained on essay-plagiarism norms, which map
poorly onto a computational replication course. Concretely, in SPPH 604
each of the following is academic misconduct:

\begin{itemize}
\tightlist
\item
  Reporting a replication result you did not actually produce, or
  submitting a repository that does not generate the numbers in your
  slides or letter.
\item
  Using the original authors' code, a public repository, a prior
  cohort's work, or course/OER code without citing it, and presenting
  the output as your own replication.
\item
  Copying another group's code, text, slides, or output --- or giving
  them yours to copy.
\item
  Copying or adapting a reviewed group's code, text, figures, output, or
  distinctive implementation into your own project; using what reviewing
  their work showed you without disclosing it; or disclosing their
  material outside the review.
\item
  Leaving your name on a submission you did not contribute to; or having
  someone outside your group write, edit, or run any part of your
  submission (see External Feedback Policy).
\item
  Undisclosed generative-AI use, or AI use outside what a given
  assessment permits (see the GenAI section below).
\item
  During the live M3 oral defense: reading from a prepared answer in
  individual Q\&A, taking help from anyone in any channel, or using any
  AI assistance or live prompting. Your slides, your own speaker notes
  and your own code are yours to use throughout --- see the GenAI
  section. Nothing in this bullet applies to a support authorized by
  your Academic Accommodation Letter --- a note-taker, live captioning,
  a scribe, or permitted assistive software --- which is arranged with
  the instructor in advance (see \emph{Academic Accommodation Letter},
  above).
\end{itemize}

By contrast, all of the following are \textbf{fine}: running the
original authors' code and saying so; copying a \texttt{svyTable1}
example from the documentation with a comment naming the source; asking
another group how they handled a \texttt{survey} package error; having
an AI write a data-wrangling function that you then audit, test, and
disclose; and asking the instructor whether any of the above is allowed.

\subsubsection{If something goes wrong}\label{if-something-goes-wrong}

Suspected academic misconduct is handled under the UBC Academic
Calendar's provisions on
\href{https://vancouver.calendar.ubc.ca/campus-wide-policies-and-regulations/student-conduct-and-discipline/discipline-academic-misconduct}{Academic
Misconduct}. That process normally begins with an opportunity to discuss
the matter with the instructor, in which any extenuating circumstances
can be raised. Where the instructor concludes that misconduct has
occurred, the instructor is required to report it to the Dean in
accordance with the applicable Faculty procedures; for graduate students
this is normally the Dean of the Faculty of Graduate and Postdoctoral
Studies or their designate. Penalties are determined through that
University process, not by the teaching team. Students are expected to
review the Student Discipline section of the
\href{https://vancouver.calendar.ubc.ca/campus-wide-policies-and-regulations/student-conduct-and-discipline/discipline-academic-misconduct}{UBC
Calendar} and to know what constitutes plagiarism and academic
misconduct.

\subsubsection{Where to get help}\label{where-to-get-help}

\begin{itemize}
\tightlist
\item
  \textbf{Ask us --- at any time.} The 48-hour guidance in the
  \emph{Communication} section applies to requests for help \emph{with
  the content of an assignment}. It does \textbf{not} apply to questions
  about academic integrity, attribution, collaboration boundaries, or
  what AI use is permitted. Those questions are always welcome, at any
  hour, by email, in lab, or in office hours. It is always cheaper to
  ask than to guess.
\item
  \textbf{UBC Academic Integrity Hub} ---
  \url{https://academicintegrity.ubc.ca/}, including its student
  guidance on what integrity means and how to avoid unintentional
  misconduct.
\item
  \textbf{UBC Learning Commons} ---
  \url{https://learningcommons.ubc.ca/} for guides on citation,
  paraphrasing, and avoiding plagiarism.
\item
  \textbf{Academic support and wellbeing resources} --- see
  \emph{University Policies}, above.
\end{itemize}

\subsubsection{External Feedback Policy}\label{external-feedback-policy}

You may seek feedback from individuals outside your group (e.g., a
supervisor, postdoctoral researcher, or subject-area expert), provided
strict adherence to the following:

\begin{itemize}
\tightlist
\item
  Approval: You must email the instructor for approval before sharing
  materials, listing the contributor's name and affiliation.
\item
  Scope: External contributors may provide comments and suggestions
  only. They must not edit text, write new material, alter code, run
  analyses, or contribute figures/tables.
\item
  Conflict of Interest: External contributors cannot be members of your
  assigned peer-review groups.
\item
  Transparency: You must note that you received external feedback in
  your submission cover note.
\end{itemize}

\subsection{Generative AI (GenAI) in Student-Submitted
Work}\label{generative-ai-genai-in-student-submitted-work}

Generative AI tools (e.g., ChatGPT, Claude, GitHub Copilot) are
\textbf{permitted and encouraged} in SPPH 604 as ``coding co-pilots''
for data wrangling and debugging. This course adopts a ``Zero to AI
Co-Pilot'' pedagogy: learning to direct, audit, and take responsibility
for AI-assisted analysis is itself a skill this course teaches. No tool
is required, none costs you anything to skip, and choosing not to use AI
will never disadvantage you.

What the course does require is that you remain the author of your
reasoning and the auditor of your output. The rules below exist to make
that workable, not to discourage you.

\subsubsection{The accountability-first
framework}\label{the-accountability-first-framework}

\begin{itemize}
\tightlist
\item
  Auditing Requirement: Students are responsible for rigorously auditing
  all AI-generated code for ``hallucinations,'' security
  vulnerabilities, and systemic health biases.
\item
  The ``No-Excuse'' Clause: Because this is a PhD-level capstone, ``The
  AI made a mistake'' is not a valid justification for errors in your
  final analysis or manuscript.
\item
  Attribution: All AI use must be disclosed in the cover note of the
  submission it went into, specifically citing which sections of code
  were AI-generated versus human-audited. Preparing for M3 is the
  exception --- it has no submission and nothing to disclose it in ---
  but AI that touched the code or slides you present at M3 is disclosed
  with M1 or M4.
\item
  Ultimate Verification: The live oral defense (M3) and the fully
  reproducible code archive (M4) serve as our primary AI-accountability
  mechanisms. An AI can write code, but you must be able to defend the
  epidemiological reasoning behind it live.
\end{itemize}

\subsubsection{Permissions by
assessment}\label{permissions-by-assessment}

Three levels are used throughout:

\begin{itemize}
\tightlist
\item
  \textcolor{blue}{\textbf{P}} --- \textbf{Permitted}, no disclosure
  required.
\item
  \textcolor{blue}{\textbf{PD}} --- \textbf{Permitted with disclosure}:
  allowed, and you record it in the AI Use Statement (format below).
\item
  \textcolor{red}{\textbf{N}} --- \textbf{Not permitted}.
\end{itemize}

{\def\LTcaptype{none} % do not increment counter
\begin{longtable}[]{@{}
  >{\raggedright\arraybackslash}p{(\linewidth - 4\tabcolsep) * \real{0.2157}}
  >{\raggedright\arraybackslash}p{(\linewidth - 4\tabcolsep) * \real{0.1176}}
  >{\raggedright\arraybackslash}p{(\linewidth - 4\tabcolsep) * \real{0.6667}}@{}}
\toprule\noalign{}
\begin{minipage}[b]{\linewidth}\raggedright
Assessment
\end{minipage} & \begin{minipage}[b]{\linewidth}\raggedright
Level
\end{minipage} & \begin{minipage}[b]{\linewidth}\raggedright
What this means here
\end{minipage} \\
\midrule\noalign{}
\endhead
\bottomrule\noalign{}
\endlastfoot
\textbf{M0} feasibility gate & \textbf{PD} & Use AI to navigate
codebooks, survey documentation, and data dictionaries. The six
eligibility criteria must be verified by you against the actual paper
and the actual data files, with real variable and file names --- an AI's
guess about what NHANES contains is not a feasibility check. \\
\textbf{L1--L6} lab exercises & \textbf{PD} & Full co-pilot use for
writing and debugging R. Flag AI-generated chunks with a one-line
comment in the script itself
(\texttt{\#\ AI-assisted:\ \textless{}tool\textgreater{}}); no separate
statement is needed for labs. Labs are in-person and graded on presence
and work done in the room, so the point is that \emph{you} leave the lab
able to do the thing. \\
\textbf{P1--P5} project increments & \textbf{PD} & Same as labs.
Disclosure is rolled up into the M1 AI Use Statement, since P1--P5 are
compiled into M1. \\
\textbf{M1} replication + diagnostics & \textbf{PD} & Full co-pilot use
for code. The \emph{interpretation} of any discrepancy between your
numbers and the published numbers must be your own --- tracing a
discrepancy to its root cause is the graded skill. \\
\textbf{M2} critique of another group's M1 & \textbf{PD}† & Use AI to
tighten language you wrote, and to fix problems on \textbf{your own
machine} while running their code --- a package that will not install,
an R version, a path their script assumes, what a generic R error means.
The rule for what you type into the tool is: \textbf{paste what is yours
or public; describe what is theirs.} Yours or public: your own
\texttt{sessionInfo()}, your own installation log, package
documentation, and an error message you have stripped of their code,
file names, variable names and numbers. Theirs: every line of their
code, and their slides, output, logs, \texttt{README}, file listing and
results --- none of it goes into any AI tool, in whole or in part,
pasted, uploaded, screenshotted or retyped. If you cannot ask the
question without quoting their work, ask the teaching team instead; that
is what we are for. You may \textbf{not} have AI generate the
substantive critique --- another group is graded on the feedback they
receive, so machine-generated criticism harms real people. Asking a
model ``what is wrong with this kind of study?'' and presenting the
generic answer is not a critique of the work in front of you, and grades
accordingly. Disclose any AI use in the cover note. \\
\textbf{M3} live oral defense & \textcolor{red}{\textbf{N}} live;
\textcolor{blue}{\textbf{P}} prep & Prepare however you like, including
rehearsing questions with an AI. Preparation is \textbf{not disclosed}:
M3 has no submission and no cover note to disclose it in. AI that
touched the code or slides you present is disclosed with M1 or M4, where
that work is submitted. \textbf{During your group's presentation}, use
everything you prepared --- your slides, your own speaker notes, your
own code and output. \textbf{During individual Q\&A, the answer must be
yours, spoken.} Look at your own slides and your own code and output as
much as you need; do not read from a prepared answer, do not take help
from anyone in any channel, and do not consult anything produced during
the session --- no AI, no live prompting, no messages, no transcription
or notetaking tool. \textbf{Approved accommodations are the exception.}
Supports authorized by your Academic Accommodation Letter --- a
note-taker, live captioning, a scribe, or permitted assistive software
--- apply at M3 as they do everywhere else in the course; you arrange
them with the instructor in advance (see \emph{Academic Accommodation
Letter}, above) and nothing has to be raised or explained during the
defense itself. M3 is the point at which the course verifies that the
reasoning is yours. \\
\textbf{M4} reanalysis letter + repo & \textbf{PD} & Full co-pilot use
for code and for drafting, editing, and polishing prose. The AI Use
Statement is a required component of the M4 submission. Note that the
Weak grade band captures work that is generic, insufficiently specific
to your data and your paper, unsupported by the output you produced,
technically incorrect, or disconnected from the estimand and its
assumptions --- unedited AI prose tends to land there on its own merits,
quite apart from disclosure. \\
\end{longtable}
}

† \emph{Two hard limits at M2: you may not have AI generate the
substantive critique, and no part of the reviewed group's work goes into
any AI tool --- paste what is yours or public, describe what is theirs.}

\subsubsection{\texorpdfstring{What you may and may not put \emph{into}
a GenAI
tool}{What you may and may not put into a GenAI tool}}\label{what-you-may-and-may-not-put-into-a-genai-tool}

\begin{itemize}
\tightlist
\item
  \textbf{Yes:} your own code, your own text, your own error messages,
  and the public de-identified study data your project uses (NHANES,
  NHIS, or comparable fully public sources). An error message that
  quotes another group's code is theirs, not yours --- strip it or
  describe the problem instead. All project data in this course is
  public and de-identified by design, so there is no
  personal-information risk in the analysis itself.
\item
  \textbf{No:} another group's unpublished slides, code, output or
  review materials, in whole or in part, pasted, uploaded, screenshotted
  or retyped (see \emph{Peer review: confidentiality and
  non-appropriation}, above, and the \textbf{M2} row of the table); and
  any personal, confidential, or restricted data, should you ever
  encounter it in adjacent work.
\item
  \textbf{Course materials} (lecture slides, lab instructions, provided
  datasets): pasting these into an AI tool for your own learning is
  permitted and is not treated as redistribution under the
  \emph{Copyright} section. Posting them publicly is governed by
  \textbf{Copyright}: the materials in the public course repository are
  CC BY 4.0 and may be posted with attribution; Canvas-only material,
  the exercise datasets and session recordings may not.
\item
  \textbf{The paper you are replicating:} uploading a published article
  to an AI tool for your own reading is permitted, subject to the
  article's own licence. Prefer open-access papers, which your M0
  eligibility criteria already require.
\end{itemize}

\subsubsection{The cover note (and your AI Use
Statement)}\label{the-cover-note-and-your-ai-use-statement}

Every submission that is not an in-lab exercise --- \textbf{M0, M1, M2,
and M4} --- is accompanied by a short \textbf{cover note}. This is the
same cover note referred to in the External Feedback Policy. It carries
three things: your per-member contribution note, any external feedback
received, and your AI Use Statement. M3 has no submission and no cover
note.

\textbf{Where it goes.} The cover note is always part of the document
you submit. Where you also submit a repository, the same text is a
plain-text or markdown file named \texttt{COVER-NOTE.md} in its root.

{\def\LTcaptype{none} % do not increment counter
\begin{longtable}[]{@{}
  >{\raggedright\arraybackslash}p{(\linewidth - 4\tabcolsep) * \real{0.1930}}
  >{\raggedright\arraybackslash}p{(\linewidth - 4\tabcolsep) * \real{0.4737}}
  >{\raggedright\arraybackslash}p{(\linewidth - 4\tabcolsep) * \real{0.3333}}@{}}
\toprule\noalign{}
\begin{minipage}[b]{\linewidth}\raggedright
Milestone
\end{minipage} & \begin{minipage}[b]{\linewidth}\raggedright
In the submitted document
\end{minipage} & \begin{minipage}[b]{\linewidth}\raggedright
In the repository
\end{minipage} \\
\midrule\noalign{}
\endhead
\bottomrule\noalign{}
\endlastfoot
\textbf{M0} & a \emph{Cover note} section at the end of the memo & ---
no repository is submitted \\
\textbf{M1} & the final slide of the deck & \texttt{COVER-NOTE.md} in
the root of the zip \\
\textbf{M2} & the final slide of the deck & --- M2 is slides only \\
\textbf{M3} & --- nothing is submitted & --- \\
\textbf{M4} & the final page of the letter & \texttt{COVER-NOTE.md} in
the root of the zip \\
\end{longtable}
}

\textbf{The AI Use Statement.} One part of the cover note, and it
carries four fields, a few lines each:

\begin{enumerate}
\def\labelenumi{\arabic{enumi}.}
\tightlist
\item
  \textbf{Tool and version} --- e.g.~``ChatGPT (GPT-5), GitHub Copilot
  in RStudio.''
\item
  \textbf{Where used} --- which tasks, and which files, scripts, or
  sections. Be specific enough that a reader could go look:
  ``\texttt{02-cohort.R}, lines 40--95 (survey design object); prose
  editing of Discussion paragraphs 2--3.''
\item
  \textbf{What you did with the output} --- how you audited it, what you
  tested, what you had to fix. This is the field the teaching team
  actually reads.
\item
  \textbf{Representative prompts} \emph{(optional but encouraged)} ---
  two or three prompts that shaped a substantive decision.
\end{enumerate}

If you used no AI at all for a submission, write one line saying so ---
that is a complete and perfectly good statement. A worked example will
be posted on Canvas alongside the M1 rubric.

\subsubsection{Consequences}\label{consequences}

Disclosed AI use is never penalized in this course. \textbf{Undisclosed
AI use, or AI use outside what the table above permits for a given
assessment, is academic misconduct} and is handled under the process
described in the \emph{Academic Integrity} section and in the UBC
Academic Calendar's
\href{https://vancouver.calendar.ubc.ca/campus-wide-policies-and-regulations/student-conduct-and-discipline/discipline-academic-misconduct}{Academic
Misconduct} provisions. Separately, and independently of any integrity
finding, unaudited AI output tends to grade in the Weak band on its own
merits.

If you are unsure whether a particular use is permitted, ask ---
including the night before a deadline. See \emph{Where to get help} in
the Academic Integrity section.

\subsection{Attendance}\label{attendance}

Tue 9:30-12 (by Zoom; 9-12 on the presentation Tuesdays in Weeks 8, 9
and 11) and Thu 10-12 (in person, SPPH 143) should be protected time for
SPPH 604, as per your registration calendar. Email the instructor in
case you can't make it to the class or lab session partially or fully
for the record. See \texttt{Late\ Assignments} section for further
details.

\subsection{Session Recordings}\label{session-recordings}

Tuesday sessions are held by Zoom and can be recorded, so that students
who miss a session can catch up (need early notice if you miss so that
instructor is aware when to record).

\begin{itemize}
\tightlist
\item
  \textbf{What is recorded, and who sees it.} Recordings are not posted
  to Canvas by default. When a recording is made, it is shared on
  request with students registered in this section this term, and is
  stored on UBC-supported services (Zoom cloud / Kaltura) reached
  through Canvas. A recording in which no student appears, is audible,
  or is otherwise identifiable --- an instructor-only walkthrough or
  demo, for example --- may be reused for teaching, including publicly.
  Any recording that includes student participation stays within this
  section this term: it is not posted publicly or shared outside the
  course without the written permission of the students concerned.
\item
  \textbf{Assessed sessions.} During your group's M1 or M2 presentation,
  and throughout your individual M3 defense, your camera must be on
  unless you have arranged an alternative with the instructor beforehand
  --- seeing you is necessary both to evaluate the presentation and to
  confirm identity. If the camera requirement is a barrier for you ---
  bandwidth, a shared or non-private space, privacy, or an accommodation
  --- contact the instructor beforehand and an alternative arrangement
  will be made.
\item
  \textbf{Recording by students.} Students may not record,
  screen-capture, or transcribe class or lab sessions, or point an AI
  notetaker or automated transcription tool at them, without the
  instructor's written permission. Students whose Academic Accommodation
  Letter includes recording are exempt, for their own study use.
\item
  No automated proctoring or invigilation software is used in this
  course.
\end{itemize}

\subsection{Learning Analytics}\label{learning-analytics}

Learning analytics includes the collection and analysis of data about
learners to improve teaching and learning. This course will be using
Canvas that capture data about student's activity and provide
information that can be used to improve the quality of teaching and
learning. In this course, the instructor plans to use analytics data to:

\begin{itemize}
\tightlist
\item
  View overall class progress
\item
  Track students' progress in order to provide them with personalized
  feedback
\item
  Review statistics on course content being accessed to support
  improvements in the course
\item
  Assess the student's participation in the course.
\end{itemize}

The instructor will invite anonymous mid-course feedback to adapt the
course in real time and ensure it meets students' learning needs and
expectations.

\subsection{Learning Resources}\label{learning-resources}

\begin{itemize}
\tightlist
\item
  \textbf{Course materials live in the public course repository} ---
  this syllabus, the lecture plans and slides, the lab handouts (L1-L6),
  the practice handouts (P1-P5), the worked model submissions, and the
  reproduction pipeline code. The repository is the authoritative source
  for all of these, and the link to it is posted on Canvas.
\item
  \textbf{Canvas carries submissions and grading}, together with the
  material this syllabus points there for: the rubric grids for M1-M4,
  the worked AI Use Statement example, Zoom links, the week-by-week
  chapter signposts, and the weekly Modules. Each week's Module is
  posted by 9 AM Monday and signposts that week's work, including which
  lab handout is in play; the handouts themselves sit in the repository
  and are available ahead of time, so you can read the lab before the
  lab.
\item
  The two required textbooks are open access and free to read online;
  they are listed under \emph{Learning Materials}, above.
\item
  Access to a computer with R software (R, RStudio, R markdown, free of
  charge) is necessary for all course work. Bring your own laptop with R
  installed to the Thursday lab in SPPH 143. Course instruction will be
  provided strictly in R.
\item
  Tuesday sessions run on Zoom. You will need a device that can run Zoom
  with a working microphone and camera, and a connection stable enough
  to present from. If any of that is a barrier, contact the instructor:
  UBC Library lends laptops, and the
  \href{https://students.ubc.ca/finances/awards-scholarships-bursaries/ubc-vancouver-technology-bursary/}{UBC
  Vancouver Technology Bursary} may have additional support.
\item
  SPPH supports for graduate students --- advising, program forms and
  checklists, funding and awards, and the SPPH Graduate Student site on
  Canvas --- are collected on the School's
  \href{https://spph.ubc.ca/current-learners/}{Current Learners} page.
\item
  Teams is used to circulate each M1 deck and code zip to its assigned
  M2 reviewers. Graded work is submitted through Canvas.
\item
  You will keep your code in a GitHub repository, but for \textbf{M1 and
  M4} you \textbf{submit a downloaded zip archive of that repository}
  (via Canvas) --- we do not require repository access.
\end{itemize}

\subsection{Copyright}\label{copyright}

All materials of this course are the intellectual property of the Course
Instructor or licensed to be used in this course by the copyright owner.
Redistribution of these materials by any means without the permission of
the copyright holder(s) constitutes a breach of copyright and may lead
to academic discipline. \textbf{Where a file came from decides whether
you may republish it.} The instructor-authored materials published in
the public course repository --- this syllabus, the lecture plans and
slides, the lab and practice handouts, the model submissions, and the
reproduction pipeline code --- are released under CC BY 4.0, as stated
in that repository's \texttt{LICENSE.md}. You may share and adapt them,
publicly and including commercially, provided you give credit, link to
the licence, and indicate what you changed. Five things sit outside that
grant and must not be redistributed: material distributed only through
Canvas and not published in the repository --- the rubric grids, the
worked AI-use example, Zoom links, and anything posted in the weekly
Modules; material licensed to the course by a third party; the exercise
datasets, which belong to their original sources and are linked rather
than copied; session recordings; and other students' work. The test is
simple: \textbf{if it is not in the public repository, do not post it.}
Session recording is covered in the \emph{Session Recordings} section
above. Posted recordings are course materials for the purposes of this
section: they are for your own academic use in this course this term and
must not be redistributed.

\begin{center}\rule{0.5\linewidth}{0.5pt}\end{center}

\textbf{Syllabus updated: September 10, 2026.}

\end{document}
